\documentclass[11pt,a4paper]{article}

\usepackage[T1]{fontenc}
\usepackage[utf8]{inputenc}
\usepackage{lmodern}
\usepackage[margin=1in]{geometry}
\usepackage{booktabs}
\usepackage{amsmath,amssymb}
\usepackage{graphicx}
\usepackage{microtype}
\usepackage{caption}
\usepackage{enumitem}
\usepackage{xcolor}
\usepackage[hidelinks]{hyperref}
\usepackage{float}
\usepackage{titlesec}

\captionsetup{font=small,labelfont=bf}
\titleformat{\section}{\normalfont\large\bfseries}{\thesection}{0.7em}{}
\titleformat{\subsection}{\normalfont\normalsize\bfseries}{\thesubsection}{0.7em}{}
\setlist{itemsep=2pt,topsep=3pt,parsep=0pt}

\newcommand{\dff}{\texttt{d\_ff}}
\newcommand{\dmodel}{\texttt{d\_model}}
\newcommand{\WO}{\ensuremath{W_O}}
\newcommand{\nats}[1]{\ensuremath{#1}\,nats}

\title{\bfseries Grafts, Depth, and Diagnostics\\[0.3em]
\large A 28M-parameter encoder under a hard parameter cap:\\
what two weeks of laptop-scale experiments established}

\author{osmansiddiquer\\\small\texttt{transformer-v1}}
\date{August 2026}

\begin{document}
\maketitle

\begin{abstract}
\noindent
We built a BERT-style encoder from scratch---attention, blocks, embeddings and MLM head, with
only the tokenizer borrowed---and pretrained it on a single 4\,GB laptop GPU under a hard
30M-parameter cap. The goal was to buy downstream accuracy with architecture and data
interventions rather than parameters. Several components worked and are reported with
controls: a staged data curriculum, replay mixing (with a corpus held at 0\% replay as the
negative control, losing \nats{0.95} against \nats{0.25} for replayed corpora), span masking,
and a two-phase procedure that installs an induction circuit into a masked language model,
raising copy-from-context accuracy from 0.8\% to 93\%. The headline goal nonetheless failed:
a \nats{0.435} (16\%) improvement in held-out MLM loss produced no measurable SST-2 gain,
because at 872 dev examples the measurement instrument (SE $\approx$ 1.2\,pp) is coarser than
every effect we produced.

The principal technical contribution is a measured correction to the folklore on
function-preserving layer grafting. Zero-initialised grafting is usually said to be exact only
in pre-LN architectures. We show it is exact in \emph{post-LN} as well, under two conditions
that are not usually stated: the block must be appended at the \emph{top} of the stack, and its
LayerNorm affine parameters must be \emph{reset} to $(\gamma{=}1,\beta{=}0)$ rather than copied
from the donor layer. Copying them costs \nats{{+}4.12}; resetting them costs
\nats{0.000} (max $|\Delta\text{logit}| = 1.3\times10^{-3}$ against a mean logit magnitude of
3.82). Inserting the same block mid-stack costs \nats{{+}0.23} regardless.

We further show that a sixth layer at this scale has no work to do under three independent
initialisations; that single-layer ablation cost over-attributes, with
$\mathrm{marginal}(L_4\mid L_3) = -1.891$; that every correlational diagnostic we built
misled at least once while causal ablation never did; that annealing corpus dominates
annealing schedule by $11\times$; and that changing only the evaluation mixture moves reported
validation loss by \nats{0.225}---over half the size of the project's entire pretraining
progress. Total cost: 35.4 GPU-hours, $\approx$3.0B tokens processed.
\end{abstract}

\tableofcontents
\vspace{1em}

%======================================================================
\section{Introduction}

The constraint that shaped this project was a hard cap of 30M parameters, on a single RTX
3050 Ti Laptop GPU with 4\,GB of VRAM. Under such a cap the interesting question is not
\emph{how well can this do}---the answer is known to be ``worse than BERT-base''---but
\emph{how much of the gap can be closed by spending the budget better}. We attacked that
question with architecture surgery (structured pruning, layer grafting, depth reallocation),
data curriculum design, masking-scheme changes, an alternative pretraining objective, and a
battery of diagnostic tooling.

This report is organised around what the experiments actually established rather than around
the chronology. Section~\ref{sec:setup} gives the setup and the evaluation protocol, which
turned out to matter more than most of the interventions. Section~\ref{sec:worked} reports the
components that demonstrably worked, with controls. Section~\ref{sec:graft} is the main
technical contribution: a measured, corrected recipe for function-preserving layer grafting in
post-LN transformers. Sections~\ref{sec:depth}--\ref{sec:width} report the negative and
diagnostic results. Section~\ref{sec:eng} collects engineering findings that cost real time.
Section~\ref{sec:scoreboard} presents the final scoreboard and the project's central negative
result. Section~\ref{sec:lessons} states what we would do differently.

Full checkpoint lineage, per-run configuration and the chronological build log are in the
repository (\texttt{checkpoints/HISTORY.md}, \texttt{WORKLOG.md}).

%======================================================================
\section{Setup}\label{sec:setup}

\subsection{Model}

\begin{table}[H]
\centering\small
\begin{tabular}{lll}
\toprule
& v1 line & v2 line (final) \\
\midrule
parameters                  & 28.733M & \textbf{27.952M} \\
layers / heads              & 4 / 4   & 6 / 4 \\
\dmodel, $d_k = d_v$        & 768, 64 & 768, 64 \\
\dff, activation            & 3072, ReLU & 1792, GELU \\
embedding                   & \multicolumn{2}{l}{ALBERT-factorised $V{\times}128$ then $128{\to}768$ (38.6M $\to$ 6.5M)} \\
MLM head                    & \multicolumn{2}{l}{tied to the word-embedding table} \\
normalisation               & \multicolumn{2}{l}{\textbf{post-LN}} \\
sequence length             & \multicolumn{2}{l}{128, contiguous packing (no padding, hence no attention mask)} \\
masking (train)             & BERT 80/10/10 scattered & geometric spans, $p{=}0.2$, clipped at 10 \\
masking (eval)              & scattered & scattered \\
\bottomrule
\end{tabular}
\caption{Architectures. The tokenizer is \texttt{allenai/OLMo-1B-hf} with an added
\texttt{[MASK]} (id 50280, vocab 50281). There is no \texttt{[CLS]}: OLMo has none and NSP was
dropped, so downstream classification mean-pools the encoder states.}
\end{table}

Parameter split of the v2 model: FFN 16.53M (59\%), embedding and projection 6.53M (23\%),
attention 4.72M (17\%), norms 0.02M. Attention holds roughly half the share it does in
BERT-base ($\approx$33\%): with $n_\text{heads}\cdot d_k = 256$ against $\dmodel = 768$, each
block writes a rank-$\leq$256 update into a 768-dimensional stream. This asymmetry recurs
below.

\subsection{Data}

Corpora were tokenised to packed \texttt{uint16} memmaps with contiguous packing (documents
joined by EOS, every block full): WikiText-103 119M tokens, UltraFineWeb-en 341M, TinyStories
475M, Cosmopedia 403M, IMDB 23M.

\subsection{Training}

bf16, AdamW, micro-batch 16 $\times$ gradient accumulation 16 (effective batch 256). Measured
VRAM fit on the 3050 Ti, which fixed the micro-batch:

\begin{center}\small
\begin{tabular}{rrrl}
\toprule
micro-batch & peak VRAM & ms/step & verdict \\
\midrule
16 & 1.84\,GB & 91   & safe default; leaves the machine usable \\
32 & 3.25\,GB & 188  & works if the GPU is otherwise idle \\
48 & 4.63\,GB & 507  & spills to WSL2 shared memory \\
64 & 6.03\,GB & 1751 & heavy spill \\
\bottomrule
\end{tabular}
\end{center}

\subsection{Evaluation protocol}\label{sec:protocol}

Every cross-model number in this report is scored on \textbf{one fixed held-out mixture}: the
tail 1\% of each Phase-A corpus, token-weighted 55/35/10 (Cosmopedia / UltraFineWeb /
TinyStories), with an identical RNG seed so that all models see byte-identical masked inputs.
Two masking schemes are reported: \texttt{scattered} (BERT 15\%, 80/10/10) and \texttt{span}
(geometric $p{=}0.2$, clipped at 10, mean span 3.80). Fifteen batches of $16\times128$; eval
noise is $\approx\pm$\nats{0.05}.

This protocol was introduced late. Adopting it retroactively converted an entirely
uncomparable set of run logs into a comparable table, and Section~\ref{sec:evalvalid} quantifies
what its absence had been costing.

%======================================================================
\section{What worked}\label{sec:worked}

\subsection{Data curriculum: the largest single lever in the project}

The staged curriculum---WikiText-103, then UltraFineWeb, then a replay mixture with
TinyStories and IMDB, then a Cosmopedia-weighted mixture---produced the two largest
improvements we recorded, both of them corpus changes rather than architecture changes.

\begin{table}[H]
\centering\small
\begin{tabular}{llrr}
\toprule
checkpoint & data added & scat\_loss & $\Delta$ \\
\midrule
\texttt{ckpt}   & WikiText-103, 119M          & 5.9866 & --- \\
\texttt{ckpt2}  & + UltraFineWeb              & 3.0389 & $-2.948$ \\
\texttt{ckpt3}  & + TinyStories + IMDB (replay) & 2.7089 & $-0.330$ \\
\texttt{ckpt\_v2} & + Cosmopedia (55\% weight), 734M & 2.3280 & $-0.381$ \\
\texttt{ckpt\_v2\_anneal\_cosmo} & 4k-step Cosmopedia anneal & \textbf{2.2738} & $-0.054$ \\
\bottomrule
\end{tabular}
\caption{Data curriculum on the fixed held-out mixture (\S\ref{sec:protocol}). The architecture
change between rows 3 and 4 is analysed separately in \S\ref{sec:width}; the corpus change is
the dominant term.}
\end{table}

Cosmopedia was weighted highest in Phase A on an explicit argument: it was the only genuinely
\emph{new} data, since \texttt{ckpt3} had already seen UltraFineWeb for $\approx$2.5 epochs,
near the point where repeats stop paying. The per-corpus breakdown below confirms this
paid---Cosmopedia improved by \nats{0.81} while everything else moved by a quarter of that.

\subsection{Replay mixing works, and we have the control}

Continued pretraining on a new corpus normally causes catastrophic forgetting. We mixed each
new phase with replay from earlier corpora. IMDB provides an unintentional but clean negative
control: it was dropped entirely from the v2 mixture while the other three corpora were
retained at 55/35/10.

\begin{table}[H]
\centering\small
\begin{tabular}{lrrrr}
\toprule
checkpoint & UltraFineWeb & TinyStories & IMDB & Cosmopedia \\
\midrule
\texttt{ckpt}  (WikiText only)          & 5.9660 & 5.4163 & 6.0371 & 6.3295 \\
\texttt{ckpt2} (+UFW)                   & 2.8882 & 2.4936 & 3.8784 & 3.4847 \\
\texttt{ckpt3} (+TinyStories+IMDB, replay) & 2.9511 & 1.0227 & 2.9255 & 3.1006 \\
\texttt{ckpt\_v2} (+Cosmopedia, replay) & 3.2132 & 1.2459 & 3.8733 & 2.2875 \\
\texttt{ckpt\_v2\_anneal\_cosmo} (final)& 3.3605 & 1.3430 & 3.9433 & 2.1129 \\
\bottomrule
\end{tabular}
\caption{Per-corpus held-out scattered MLM loss. Each column is a separate held-out split;
columns are not comparable to each other, only down.}
\end{table}

Reading the \texttt{ckpt3} $\to$ \texttt{ckpt\_v2} transition against the replay weight:

\begin{center}\small
\begin{tabular}{lrr}
\toprule
corpus & weight in the v2 mixture & $\Delta$ loss \\
\midrule
Cosmopedia    & 55\% (new)       & $-0.813$ \\
UltraFineWeb  & 35\%             & $+0.262$ \\
TinyStories   & 10\%             & $+0.223$ \\
IMDB          & \textbf{0\%}     & \textbf{+0.948} \\
\bottomrule
\end{tabular}
\end{center}

Any replay at all held forgetting to $\approx$\nats{0.25}; zero replay cost \nats{0.95}, close
to four times as much, and dropping from 35\% to 10\% cost nothing measurable. The practical
reading is that the replay \emph{floor} matters much more than the replay \emph{fraction}: a
token budget of 10\% is enough, and 0\% is not.

The same table also characterises what annealing does, honestly. The Cosmopedia anneal improved
Cosmopedia by \nats{0.175} and degraded all three other corpora (UltraFineWeb $+0.147$,
TinyStories $+0.098$, IMDB $+0.070$). Annealing \emph{specialises}; it does not uniformly
improve. This is why the anneal simultaneously measured $-0.054$ on the Phase-A mixture and
$+0.057$ on IMDB, an apparent contradiction at the time.

\subsection{Span masking}

Following SpanBERT, the v2 line trains on geometric span masking ($p{=}0.2$, clipped at 10,
mean span 3.80) and evaluates on scattered masking so the logged accuracy stays comparable with
earlier runs. Span masking removes the ability to infill from immediate neighbours, forcing
representation rather than local autocomplete.

It worked, on two independent readings. First, the model improved \emph{more} on the harder
task than on the easier one: from \texttt{ckpt3} to the final model, scattered accuracy rose
4.6\,pp while span accuracy rose 6.7\,pp. Second, the intrinsic difficulty gap between the two
schemes narrowed, from \nats{{+}1.932} for the scattered-trained \texttt{ckpt3} to
\nats{{+}1.640} for the span-trained final model---a \nats{0.29} narrowing that is six times
eval noise.

That gap is large enough to be worth stating on its own, because it is a frequent source of
false alarms. Measured on one fixed model with identical data and seed:

\begin{center}\small
\begin{tabular}{lrrrr}
\toprule
masking scheme (\texttt{mlm\_prob} 0.15) & mean span & mlm\_loss & masked\_acc & vs scattered \\
\midrule
scattered            & 1.16 & 2.3655 & 0.5814 & --- \\
uniform 2--4         & 2.94 & 3.5002 & 0.4362 & $+1.135$ \\
geometric $p{=}0.2$  & 3.64 & 4.1054 & 0.3662 & $\mathbf{+1.740}$ \\
geometric $p{=}0.1$  & 5.36 & 4.4758 & 0.3186 & $+2.110$ \\
\bottomrule
\end{tabular}
\end{center}

A model that trains on spans and evaluates on scattered will show train and validation curves
$\approx$\nats{1.7} apart \emph{by design}. Anyone reading those logs without this table would
diagnose a catastrophic bug.

\subsection{WSD with the decay phase left unspent}

\texttt{ckpt3} used cosine decay to zero and demonstrated that it bought nothing: the best
validation loss arrived at step 6{,}400 while the learning rate was still at 55\% of peak, and
the fully annealed endpoint was \emph{worse} than the mid-run best. Every later run therefore
used Warmup--Stable--Decay with \texttt{decay\_frac = 0}, ending flat.

Two benefits followed. The run stays extendable without a re-warm, and the decay phase becomes
a separate, controllable experiment---which is exactly what made the corpus-controlled
annealing comparison of \S\ref{sec:anneal} possible at all. Spending the decay is a decision;
it should not be a side effect of the schedule.

\subsection{Installing an induction circuit into a masked language model}\label{sec:autoreg}

We coerced the encoder into an autoregressive generator---append three \texttt{[MASK]} tokens,
supervise and commit slot 1, slide the window---as a probe of what the MLM had absorbed. Two
capabilities were missing and both were installable in minutes of CPU-class fine-tuning.

\begin{itemize}
\item \textbf{Termination.} The base model never emits EOS in a generation loop. Oversampling
EOS positions (30\% of supervised slots) during a 400-step continuation fine-tune fixes this.
\item \textbf{Copy-from-context is absent at baseline: 0.8\%.} Induction requires a two-layer
circuit---a previous-token head feeding a copy head---and four heads per layer leave no room
for one to form incidentally under a masked objective. A 600-step fine-tune on two-entity
templates (25\% of the batch, the rest replay) raises copy accuracy to \textbf{93.3\%} on
trained-style prompts and \textbf{80.0\%} on 200 entities held out of the training pool, in
6.5 minutes.
\end{itemize}

The generalisation to held-out entities is what makes this a circuit rather than memorisation.
The result is that the capability was \emph{cheap and absent}: nothing in 2.1B tokens of MLM
pretraining gave the model a reason to build it, and once given a reason it built it in
6.5 minutes.

Two honest caveats. The fine-tune over-corrects at high copy weight---\emph{``We went to Paris
last year with Maria. I loved chatting with''} completes to \emph{``Paris last year with
Maria''}. And there is an open failure: generation loops on prompts shorter than 8 tokens,
because training sampled context length uniformly over $[8, 125]$ while generation
\emph{starts} at the prompt length. All four sub-8-token prompts break; all five longer prompts
are clean. The fix (log-uniform context sampling from length 1) is identified but not
implemented.

\subsection{The diagnostic tooling}

The tools in \texttt{tools/} are themselves a result. Three earned their cost repeatedly:
\texttt{span\_eval.py} (the fixed-protocol evaluation of \S\ref{sec:protocol}),
\texttt{layer\_contrib.py} (causal ablation cost, single and pairwise, \S\ref{sec:depth}), and
\texttt{audit\_eval.py} (which mechanically checks the failure modes of
\S\ref{sec:evalvalid}). Each converted a recurring argument into a number.

%======================================================================
\section{Function-preserving grafts in post-LN transformers}\label{sec:graft}

\subsection{The problem}

Growing a trained network by adding layers---progressive stacking, bert2BERT, Net2Net---is
meant to be near-free: the new layer starts from a useful function and training continues from
a good point. The standard trick for making it \emph{exactly} free is to zero the new block's
output projections so that the block is an identity at step 0.

This is folklore for \textbf{pre-LN} architectures, where a block computes $x + f(\mathrm{LN}(x))$
and $f \equiv 0$ genuinely is the identity. Our model is \textbf{post-LN}:
\begin{align}
h   &= \mathrm{LN}_1\!\left(\mathrm{MHA}(x) + x\right), \\
\text{out} &= \mathrm{LN}_2\!\left(\mathrm{FFN}(h) + h\right).
\end{align}
With $\WO = 0$ and $W_\text{out} = 0$ this collapses to
$\text{out} = \mathrm{LN}_2(\mathrm{LN}_1(x))$---two extra LayerNorms, not an identity. We
measured what that costs and found the folklore is recoverable in post-LN, subject to two
conditions that are not usually stated.

\subsection{Measurement}

All variants below are built from the same trained 4-layer checkpoint (\texttt{ckpt3}) by
appending \emph{one} block, and are scored under the protocol of \S\ref{sec:protocol}. Because
building a v2-package model also swaps ReLU for GELU, $\Delta$ is reported against the
GELU-only reference so the activation change is factored out.

\begin{table}[H]
\centering\small
\begin{tabular}{lrrr}
\toprule
variant (one block appended to a trained 4L stack) & scat\_loss & $\Delta$ vs ref & $\max|\Delta\text{logit}|$ \\
\midrule
\texttt{ckpt3}, 4L, ReLU --- the source model                       & 2.7089 & ---      & --- \\
GELU swap only, still 4L --- \textbf{the reference}                 & 2.9776 & $0.000$  & --- \\
$\WO{=}0$, $W_\text{out}{=}0$, LN affines reset to $(1,0)$          & \textbf{2.9776} & $\mathbf{+0.000}$ & $\mathbf{0.0013}$ \\
$\WO{=}0$, $W_\text{out}{=}0$, LN affines \emph{copied from donor}  & 7.1016 & $+4.124$ & $37.7$ \\
only $W_\text{out}{=}0$ (attention live), LN copied                 & 7.7290 & $+4.751$ & --- \\
full copy-stack --- \emph{what this project actually did}           & 8.5488 & $+5.571$ & --- \\
\bottomrule
\end{tabular}
\caption{Decomposing the cost of adding a layer to a trained post-LN stack.}
\label{tab:graft}
\end{table}

The mean logit magnitude is 3.82, so $\max|\Delta\text{logit}| = 1.3\times10^{-3}$ is float32
round-off (a relative error of $3.3\times10^{-4}$): the corrected graft is \textbf{exactly} function-preserving. The naive graft
destroys the model.

\subsection{Why it works, and where it stops working}

A zero-output post-LN block computes $\mathrm{LN}(\mathrm{LN}(x))$, which with unit affines
equals $\mathrm{LN}(x)$---a per-token standardisation of $x$. That is \emph{not} $x$. But
LayerNorm is invariant to per-token affine rescaling,
\begin{equation}
\mathrm{LN}(a\,x + b\,\mathbf{1}) = \mathrm{LN}(x), \qquad a > 0,\ b \in \mathbb{R},
\label{eq:lninv}
\end{equation}
and $\mathrm{LN}(x)$ is itself exactly such a rescaling of $x$ (with $a = 1/\sigma$,
$b = -\mu/\sigma$). So if the \emph{next} consumer of the stream is another LayerNorm, the
standardisation is annihilated and the graft is invisible.

At the \textbf{top} of the stack the only consumer is the encoder's final LayerNorm, and
(\ref{eq:lninv}) applies exactly. \textbf{Mid-stack}, the next block's attention reads the
residual stream \emph{directly}---attention logits scale as the square of the stream
magnitude---while the graft has just stripped the previous block's learned output scale
$(\gamma,\beta)$. Equation~(\ref{eq:lninv}) no longer applies:

\begin{center}\small
\begin{tabular}{lrr}
\toprule
insertion point of the zero-graft block (5L/3072) & scat\_loss & $\Delta$ \\
\midrule
nowhere --- 4L reference                    & 2.9776 & --- \\
position 0 (before all trained blocks)      & 3.0646 & $+0.087$ \\
position 2 (mid-stack)                      & 3.2021 & $+0.225$ \\
position 4 (appended at top)                & \textbf{2.9776} & $\mathbf{+0.000}$ \\
\bottomrule
\end{tabular}
\end{center}

Position 0 is cheaper than position 2 only because the embedding-plus-positional stream it
re-standardises is less shaped by a learned LayerNorm than a mid-stack activation is.

\subsection{The counter-intuitive term}

The dominant error is \textbf{copying the donor block's LayerNorm affines}: \nats{{+}4.12} of
the \nats{{+}5.57} total. Every instinct in progressive stacking says to copy everything from a
trained layer---that is the entire argument for stacking over random initialisation. For the
LayerNorm affines specifically, copying is worse than resetting by more than four nats, because
$\gamma$ and $\beta$ encode \emph{what scale the next layer expects}, and a newly inserted
block inherits the wrong expectation.

\subsection{Practical rules}

\begin{enumerate}[label=(\alph*)]
\item \textbf{Post-LN:} zero-init grafting is exact \emph{only at the top of the stack}, and
only with $\WO = 0$, $W_\text{out} = 0$, \textbf{and LayerNorm affines reset to $(1,0)$}.
Mid-stack, budget $\approx$\nats{0.2}.
\item \textbf{Copy weights, reset norms.} Inherit the residual-branch parameters; do not
inherit $\gamma$ and $\beta$.
\item \textbf{Verify with logits, not loss.} Our first check compared loss on a small batch and
read a plausible-looking delta. Comparing $\max|\Delta\text{logit}|$ against the mean logit
magnitude turns a judgement call into a pass/fail: $1.3\times10^{-3}$ versus $37.7$ is not a
debate. The verification belongs \emph{inside} the tool that performs the graft.
\item \textbf{A destroyed output does not mean destroyed knowledge.} \texttt{ckpt\_v2\_init}
scored \nats{8.26}. It recovered to \nats{2.71}---its parent's level---within
1{,}000 steps, and to \nats{2.38} by step 14{,}200, a level \texttt{ckpt3} needed $\approx$42{,}000
cumulative steps to reach. Loss immediately after a graft is a very poor proxy for how much
transferred.
\end{enumerate}

%======================================================================
\section{Depth at 28M: the sixth layer}\label{sec:depth}

We grew 4 layers to 6 by cutting \dff{} from 3072 to 1792 and spending the savings on depth.
Layer 5 pulled its weight. \textbf{Layer 4 did not.}

\subsection{Ablation cost: a causal metric}

Correlational metrics (utilisation, effective rank, head similarity) describe what a layer
\emph{looks like}. We wanted what it is \emph{worth}, so we defined \textbf{ablation cost}: zero
the layer's output projections ($\WO$, $l_2$, $b_2$), leaving the residual path intact, and
measure the increase in held-out loss.

\begin{center}\small
\begin{tabular}{lrrrrrr}
\toprule
layer & 0 & 1 & 2 & 3 & \textbf{4} & 5 \\
\midrule
ablation cost (nats) & 1.186 & 1.683 & 1.749 & 1.594 & \textbf{0.073} & 0.963 \\
\bottomrule
\end{tabular}
\end{center}

Layer 4 is worth 7.6\% of what its neighbour layer 5 is worth, and 1.0\% of the stack's total
ablation cost, while holding an identical 12.7\% share of the parameters.

\subsection{Three initialisations, one answer}

\begin{center}\small
\begin{tabular}{lrr}
\toprule
initialisation of layer 4 & steps trained & final ablation cost \\
\midrule
stacked copy of a trained layer                     & 9{,}600  & $\approx 0$ \\
zero-init graft, $5\times$ LR, no weight decay      & 4{,}600  & $\approx 0$ \\
standard $\mathcal{N}(0,0.02)$, \emph{rest of the network frozen} & 18{,}800 & $\mathbf{+0.062}$ \\
\bottomrule
\end{tabular}
\end{center}

The third run is the decisive one. Freezing embeddings, the MLM head and blocks
$0,1,2,3,5$, leaving 3.545M of 27.952M parameters trainable (12.7\%), removes every escape
route: the rest of the network \emph{cannot} route around layer 4, so whatever it can learn, it
must. It climbed from 0 to \nats{{+}0.062} over 18{,}800 steps and flattened, which is
$\approx$6\% of layer 5's contribution.

\textbf{Verdict: position, not initialisation.} At this depth and scale there is no work
available at slot 4 that existing redundancy in layers 1, 2 and 5 does not already absorb.
Transplanting the frozen-trained layer back into the joint run improved loss by
\nats{-0.056} instantly---and joint training gave most of it back within 200 steps. The joint
optimiser actively \emph{prefers} the configuration in which layer 4 does nothing.

\subsection{Pairwise marginals change how ablation should be read}

Single-layer ablation cost is not a clean attribution. We measured the full pairwise matrix,
with $\mathrm{marginal}(i \mid j) = \mathrm{cost}(\{i,j\}) - \mathrm{cost}(\{j\})$:

\begin{center}\small
\begin{tabular}{lrrrrr}
\toprule
$j =$ & 0 & 1 & 2 & \textbf{3} & 5 \\
\midrule
$\mathrm{marginal}(4 \mid j)$ & $-0.012$ & $+0.031$ & $-0.019$ & $\mathbf{-1.891}$ & $+0.005$ \\
\bottomrule
\end{tabular}
\end{center}

Layer 4 is worth $\approx$0 given \emph{any} other layer removed, so redundancy is not the
explanation---nothing was covering for it. But at $j = 3$, \textbf{removing the useless layer 4
halves the apparent cost of removing layer 3} ($3.437 \to 1.545$).

The mechanism is fragility amplification. Ablating layer 3 produces an off-distribution
residual stream; layer 4, trained only on in-distribution input, amplifies that corruption
rather than absorbing it. Layer 3's headline cost of \nats{3.44} is therefore
$\approx$55\% damage routed through layer 4 and only $\approx$45\% its own contribution.

\textbf{Rule: single-layer ablation cost over-attributes to layers sitting below a fragile
downstream layer.} For pruning decisions the pairwise marginal is the number wanted, and it can
differ from the single by a factor of two.

%======================================================================
\section{Diagnostics that mislead}\label{sec:diag}

We built five diagnostics. Four of them lied at least once, in ways worth cataloguing.

\subsection{Head similarity had a null of $\approx$0.97, not 0}

The maximum pairwise cosine between heads' attention matrices read 0.98 on layer 3, apparently
indicating duplicate heads. It does not: \textbf{an untrained model reads 0.95--0.99 as well.}
Attention rows are probability distributions over 128 positions; an undifferentiated head sits
near-uniform, and the shared $1/T$ floor dominates the cosine. The metric could not distinguish
``duplicated heads'' from ``untrained heads''---precisely the two cases it existed to separate.

Centring across heads before taking the cosine restores a meaningful null. Measured on a
freshly initialised model: \textbf{raw 0.95--0.99, centred $-0.23$ to $-0.07$}. The null is
\emph{negative} because four centred vectors sum to zero, so undifferentiated heads are
mutually anti-correlated.

A second bug compounded the first: zeroing the diagonal before the argmax both won the argmax
and clamped the reported maximum, so the tool reported the nonsense pair $h_0 \sim h_0$ at
0.00. The diagonal must be $-\infty$.

After the fix, layer 3's 0.98 survived (genuinely duplicated) while layers 4 and 5's 0.86 and
0.98 collapsed to 0.18 and 0.27---they were untrained, not duplicated. Every conclusion drawn
from the uncentred metric had to be discarded.

\subsection{The same failure one level up: latent-target pretraining}

We ran a data2vec/JEPA-style latent-prediction branch (EMA teacher, stop-gradient, narrow
predictor, momentum $0.996 \to 0.999$) for 15{,}258 steps on 500M tokens. Every internal metric
was encouraging.

\begin{center}\small
\begin{tabular}{rrrrrr}
\toprule
step & centred cos & raw cos & constant-predictor baseline & eff.\ rank & target std \\
\midrule
200     & 0.0357 & 0.886 & 0.916 & 240.9 & 0.349 \\
8{,}600 & 0.6161 & 0.944 & 0.907 & 232.9 & 0.363 \\
15{,}200& \textbf{0.6256} & 0.940 & 0.895 & 234.8 & 0.382 \\
\bottomrule
\end{tabular}
\end{center}

Note the raw cosine at step 200: \textbf{0.886, against a constant-predictor baseline of
0.916}. A predictor that ignored its input entirely and emitted the dataset mean would have
scored \emph{higher}. Representations are anisotropic; almost any two vectors are similar. The
honest number was 0.0357.

We caught that one. What we did not catch until downstream evaluation is that the \emph{whole
class} of internal metric was uninformative here: centred cosine rose $17\times$, effective
rank held steady, target standard deviation improved monotonically---and the downstream result
was indistinguishable from the far simpler MLM line (\S\ref{sec:scoreboard}).

\subsection{Zero-output layers poison run-level aggregates}

FFN-slack and attention-health tools reported \texttt{slack = 100\%}, \texttt{util = 1.00},
\texttt{out\_rank = nan} for a freshly grafted layer whose output projection is identically
zero. Those rows entered the run-level means and flipped an automated verdict to ``OK'' for
several monitoring cycles. Any health metric of the form \emph{``how much of this layer's
output matters''} is undefined when the output is zero by construction; the tool must detect
and exclude that case rather than average it in.

\subsection{The general lesson}

\textbf{Every proxy metric needs its null measured on an untrained model of the same
architecture, and its degenerate cases enumerated, before it is trusted once.} Ablation cost
was the only diagnostic we never had to retract, because it is defined by an intervention on
the model rather than by a statistic of its activations. It is also $\approx$40$\times$ more
expensive per reading. That trade was worth it.

%======================================================================
\section{Annealing: corpus dominates schedule}\label{sec:anneal}

Phase A used WSD with \texttt{decay\_frac = 0}, deliberately leaving the decay phase unspent. We
then spent it twice from the \emph{identical} checkpoint with \emph{identical} learning rate,
schedule, seed and step count (cosine $2\times10^{-4} \to 0$, no warmup, 4{,}000 steps),
varying only the corpus.

\begin{center}\small
\begin{tabular}{lrll}
\toprule
anneal corpus & own val\_loss & held-out result vs parent & outcome \\
\midrule
TinyStories & \textbf{0.9989} & fell \emph{below} \texttt{ckpt3} within 1{,}800 steps & abandoned \\
Cosmopedia  & 2.0487 & $-0.054$ nats, $+1.18$\,pp span acc & kept; final model \\
\bottomrule
\end{tabular}
\end{center}

TinyStories produced by far the lowest validation loss the project ever recorded, and it was
pure artefact: TinyStories is deliberately constructed with a $\approx$1{,}500-word vocabulary
at a three-to-four-year-old reading level. Annealing on it \emph{narrows the lexicon}; it does
not teach fluency. It lost three quarters of v2's entire advantage over \texttt{ckpt3} in 800
steps. Cosmopedia---the same schedule on synthetic textbook prose---did \textbf{11$\times$ less
damage} and net helped.

Because everything except the corpus was held fixed, this isolates data register as the causal
variable. \textbf{Annealing is a data-selection decision, not a schedule trick, and the loss on
the anneal corpus itself carries no information about whether it worked.}

%======================================================================
\section{Evaluation validity}\label{sec:evalvalid}

Late in the project a reported validation loss of $\approx$2.04 looked implausibly good---a
well-trained 124M-parameter GPT-2 on 10B tokens lands near 3.3, and we would have been beating
it with a quarter of the parameters and 7\% of the data. We audited it against the standard
failure modes.

\paragraph{1. It was a different task.} Masked language modelling is not autoregressive language
modelling; bidirectional context makes the token far easier to predict. There was no valid
comparison to GPT-2 in the first place. A category error rather than a bug, but the most
consequential item on the list.

\paragraph{2. The eval mixture, not the model, was doing most of the work.} For the \emph{same}
checkpoint:

\begin{center}\small
\begin{tabular}{lr}
\toprule
evaluation mixture & scat\_loss \\
\midrule
Cosmopedia only (what that run logged)   & 2.0487 \\
Phase-A mixture, 55/35/10 (the protocol) & 2.2738 \\
\bottomrule
\end{tabular}
\end{center}

\nats{0.225} from the evaluation set alone---over half the size of the entire project's
pretraining improvement (\nats{0.435}). Every run in this repository evaluates on its own
training mixture by default, which makes the \texttt{val\_loss} column of any cross-run table
meaningless.

\paragraph{3. Cross-run contamination is real and easy to miss.} The split is contiguous
(validation is the tail 1\%), so it is clean \emph{within} a run. But \texttt{limit\_tokens}
changes where the boundary falls, and weights are inherited across runs:

\begin{center}\small\ttfamily
\begin{tabular}{ll}
tinystories \quad n = 73.4M & train=[0, 72.6M) \quad val=[72.6M, 73.4M) \\
\quad prior run & train=[0, 165.3M) $\to$ \textbf{VAL INSIDE PRIOR TRAIN: CONTAMINATED} \\
\end{tabular}
\end{center}

v2's TinyStories validation range sits entirely inside \texttt{ckpt3}'s training range, and v2
inherited \texttt{ckpt3}'s weights. UltraFineWeb was clean; TinyStories was not.

\paragraph{4. Averaging bias was real but negligible.} Mean-of-batch-means 2.2743 against
token-weighted 2.2760: a $-0.0016$ nat bias. Worth checking once; not worth worrying about.

\paragraph{5. Data was not pathologically repetitive.} 0 of 320 validation blocks byte-matched
any of 199{,}989 sampled training blocks. The top-8 scored tokens account for 22.5\% of scored
positions across 2{,}357 distinct types---normal for English.

\paragraph{6. The generalisation gap is small but nonzero.} Train-range 2.1427 against held-out
2.3251: \nats{{+}0.18} and $-1.2$ accuracy points. Consistent with a capacity-bound,
zero-dropout model that is not overfitting.

%======================================================================
\section{Width--depth reallocation}\label{sec:width}

Three measurements on \texttt{ckpt3} motivated the v2 architecture:

\begin{itemize}
\item \textbf{\dff{} was $\approx$3.5$\times$ oversized.} 99\% of FFN activation variance fit in
133--866 of 3{,}072 directions, depending on layer.
\item \textbf{Layer 0 had 2{,}697 of 3{,}072 units that never fired once} over 2.0M tokens under
ReLU.
\item \textbf{Residual-stream effective rank was still climbing at the last layer}
($149 \to 269 \to 483 \to 523$ of 768), so depth had headroom where width did not.
\end{itemize}

We therefore cut \dff{} to 1792, swapped ReLU for GELU---and, importantly, scored the pruning
importances \emph{under GELU}, because scoring under ReLU would discard units whose
pre-activations are negative (2{,}755 of them in layer 0) that GELU passes with a small
non-zero output---and spent the savings on two extra layers, landing at 27.952M under the cap.

\textbf{The width cut cost almost nothing}, as predicted: structured pruning $3072 \to 1792$
cost \nats{{+}0.083} on top of the activation swap. After training, $r_{99}/\dff \leq 0.56$ on
every layer---width was still not the binding constraint even at 1792.

\textbf{The depth it funded returned almost nothing}: layer 5 pulled its weight, layer 4 did not
(\S\ref{sec:depth}), and the end-to-end gain did not survive to the downstream task
(\S\ref{sec:scoreboard}).

This is what scaling laws predict---performance depends on parameters, data and compute, and
only weakly on architectural shape. We ran the experiment carefully enough to confirm rather
than assume it, but the honest summary is that \textbf{reallocating a fixed parameter budget
between width and depth is a rounding error compared with changing the budget.}

%======================================================================
\section{Engineering results}\label{sec:eng}

\paragraph{Freezing 87\% of parameters buys 1.6$\times$, not 8$\times$.} The frozen-stack
experiment trained only layer 4---12.7\% of parameters---at 0.858\,s/step against 1.397\,s/step
for full joint training: $1.63\times$. The reason is FLOP structure, not parameter count. The
forward pass is unchanged (all six layers), and the backward pass must still traverse layer 5
to reach layer 4; only the four layers \emph{below} the trainable one skip backward. Predicted
$1.8\times$, measured $1.63\times$. \textbf{Count the backward passes you actually skip, not the
parameters you froze.}

\paragraph{Budget by steps, not wall-clock, on a machine that sleeps.} A
\texttt{--max-seconds 10800} budget using \texttt{time.time()} counted sleep and terminated a
run at step 1{,}851 of 8{,}800 after 37 minutes of real compute.

\paragraph{A watchdog inside the thing it watches is not a watchdog.} Every recovery mechanism
we built initially lived inside the session that sleep suspended, dying at the same instant as
the training it existed to observe. Moving it to cron covered crashes but not Windows sleep,
which \emph{destroys} the WSL2 VM---confirmed by two distinct boot IDs with an 8h38m hole in
\texttt{journalctl --list-boots}. There is no kernel running to fire cron.

\paragraph{A watchdog that spawns a daemon can disable itself.} The supervisor held its lock on
file descriptor 9; \texttt{setsid} children inherit open descriptors, so the long-lived
dashboard \emph{it} spawned captured fd 9 and held the lock for its lifetime. Every subsequent
cron tick died at \texttt{flock -n 9} and exited 0, silently, for $\approx$5 hours. The one
recovery test that passed did so only because the dashboard alive at that moment had been
started from a shell. \textbf{The act of the supervisor launching a daemon is what disabled the
supervisor---that test could not have caught the bug however carefully it was run.} Verifying a
watchdog once tells you nothing about whether it stays armed.

\paragraph{\texttt{\$\{VAR:-default\}} substitutes on empty, not merely on unset.}
\texttt{BOOST\_LAYERS=""} fell through to the default and silently launched two runs with a
$5\times$ learning rate that had been explicitly ruled out. \texttt{\$\{VAR-default\}} is the
correct form. Cost: two killed runs and two checkpoint restores.

\paragraph{Grep for a process name and you will match your own shell.} A GPU-wait loop using
\texttt{ps | awk '/pretrain\_v2/'} matched any command that merely \emph{mentioned} the string,
including the monitoring commands, and wedged forever. Poll the resource
(\texttt{nvidia-smi} VRAM), not the name.

\paragraph{Sparse-file apparent size is not disk usage.} Packed \texttt{.bin} memmaps are
allocated at target size and written incrementally; the holes cost zero real blocks.
Truncating them to \texttt{tokens\_written} reduced apparent size from 3.7\,GB to 2.4\,GB and
freed nothing---and requires updating the manifest in lockstep, since the dataset memmaps with
\texttt{shape=(target\_tokens,)}.

%======================================================================
\section{Scoreboard, and the central negative result}\label{sec:scoreboard}

\subsection{Pretraining progressed clearly}

\begin{table}[H]
\centering\footnotesize\setlength{\tabcolsep}{4.5pt}
\begin{tabular}{llrrrr}
\toprule
checkpoint & step / arch & scat\_loss & scat\_acc & span\_loss & span\_acc \\
\midrule
\texttt{ckpt} --- WikiText-103            & 8{,}800 \; 4L/3072  & 5.9866 & 0.2531 & 6.9932 & 0.1695 \\
\texttt{ckpt2} --- +UltraFineWeb          & 20{,}789 \; 4L/3072 & 3.0389 & 0.4819 & 4.9318 & 0.2884 \\
\texttt{ckpt3} --- +TinyStories+IMDB      & 12{,}964 \; 4L/3072 & 2.7089 & 0.5295 & 4.6408 & 0.3118 \\
\texttt{ckpt\_v2\_init} --- graft, untrained & 0 \; 6L/1792     & 8.2635 & 0.0125 & 8.4851 & 0.0075 \\
\texttt{ckpt\_v2} @ 14.2k                 & 14{,}200 \; 6L/1792 & 2.3828 & 0.5661 & 4.0408 & 0.3588 \\
\texttt{ckpt\_v2} @ 18.4k (pre-anneal)    & 18{,}400 \; 6L/1792 & 2.3280 & 0.5759 & 3.9809 & 0.3671 \\
\texttt{ckpt\_v2\_l4only} --- frozen-stack L4 & 18{,}800 \; 6L/1792 & 2.3136 & 0.5781 & 3.9658 & 0.3695 \\
\textbf{\texttt{ckpt\_v2\_anneal\_cosmo}} --- final & +4{,}000 \; 6L/1792 & \textbf{2.2738} & 0.5759 & \textbf{3.9135} & \textbf{0.3789} \\
\bottomrule
\end{tabular}
\caption{All rows scored under the fixed protocol of \S\ref{sec:protocol}. \texttt{ckpt3} to
final is $-0.435$ nats scattered, $-0.727$ nats span, $+6.7$\,pp span accuracy---$\approx9\times$
eval noise and unambiguously real. The latent-target branch is absent because it never trained
an MLM head; scoring it under MLM reads \nats{11.98} and means nothing.}
\end{table}

\subsection{Downstream did not}

SST-2 dev is 872 examples: SE $\approx$ 1.2\,pp. STS-B dev is 1500: SE $\approx$ 1.7\,pp.

\begin{table}[H]
\centering\footnotesize\setlength{\tabcolsep}{5pt}
\begin{tabular}{lrrrrrr}
\toprule
model & \multicolumn{1}{c}{SST-2} & SST-2 & SST-2 & STS-B & STS-B & STS-B \\
      & \multicolumn{1}{c}{standalone} & probe & ft & probe & ft & 0-shot \\
\midrule
\texttt{ckpt3}                    & 0.8567 & 0.7741 & 0.8509 & 0.5971 & 0.4199 & 0.5610 \\
\texttt{ckpt\_jepa}               & 0.8429 & 0.7362 & \textbf{0.8647} & 0.6035 & \textbf{0.6111} & 0.4600 \\
\texttt{ckpt\_v2\_anneal\_cosmo}  & \textbf{0.8601} & 0.7431 & 0.8509 & 0.5853 & 0.1434\,$^\dagger$ & 0.5558 \\
\emph{BERT-base (110M)}           & \emph{$\approx$0.927} & & & & & \\
\bottomrule
\end{tabular}
\end{table}

\noindent
\textbf{This is the project's central negative result.} A \nats{-0.435} improvement in held-out
MLM loss produced a \textbf{0.0\,pp} change on the SST-2 fine-tune harness and $+0.34$\,pp
(0.3\,SE) on the standalone harness. Nothing here is distinguishable from noise: the full
spread across three very different pretraining regimes is 0.8509--0.8647, i.e.\ \textbf{1.1\,SE}.

Worse for interpretation, \textbf{the two SST-2 harnesses disagree by more than any of our
interventions moved the number.} The same final checkpoint reads 0.8601 standalone and 0.8509
in the benchmark---a 0.92\,pp gap arising from selection protocol and split handling alone,
against a 0.34\,pp effect size. And the latent-target branch, which the standalone harness
ranked last, is the \emph{best} model on both benchmark fine-tunes.

The honest reading is that at 28M parameters on an 872-example dev set, \textbf{we built a
measurement instrument too coarse to detect what we were optimising.} Resolving a 1\,pp
difference at this standard error would need roughly four times more labelled evaluation data.
That should have been checked before the third architecture experiment, not after the last one.

The remaining $\approx$7\,pp gap to BERT-base is capacity, not allocation. The final model's
lineage processed $\approx$2.18B tokens---78 tokens per parameter, far past any compute-optimal
ratio---while validation loss plateaued at the eval-noise floor and zero dropout produced no
overfitting. All three signals point the same way.

\vspace{0.5em}
\noindent
$^\dagger$ \texttt{stsb finetune = 0.1434} is broken, not a result: a fine-tune cannot
legitimately land four times below its own frozen probe (0.5853). It read 0.29--0.32 throughout
training, i.e.\ it diverged. The identical code gave 0.4199 and 0.6111 on other backbones.
Undiagnosed.

%======================================================================
\section{What we would do differently}\label{sec:lessons}

\begin{enumerate}
\item \textbf{Size the evaluation before sizing the experiment.} SE $\approx$ 1.2\,pp on SST-2
dev makes any intervention worth less than $\approx$2.5\,pp undetectable. We ran three
architecture experiments against an instrument that could not resolve their effects. Either
pick a task with a larger dev set, average over seeds, or accept up front that the experiment
is qualitative.
\item \textbf{Measure every proxy metric's null on an untrained model before trusting it once.}
The head-similarity metric read 0.97 at random initialisation. Ten minutes of work would have
caught it; instead it invalidated a week of readings.
\item \textbf{Prefer causal metrics.} Ablation cost was $\approx$40$\times$ more expensive per
reading and the only diagnostic never retracted. Budget for it.
\item \textbf{Use pairwise marginals, not singles, for pruning decisions.}
$\mathrm{marginal}(4 \mid 3) = -1.891$ means single-layer costs can be twice wrong, in the
direction that matters.
\item \textbf{Fix the evaluation mixture globally on day one.} Per-run \texttt{val\_loss} on a
per-run mixture is not a metric; it is a diary entry. A single frozen held-out set turned an
uncomparable results table into a comparable one, retroactively, in one afternoon.
\item \textbf{Keep a replay floor, not a replay fraction.} 10\% replay retained as well as 35\%;
0\% cost four times as much.
\item \textbf{Do not spend a parameter budget reallocating shape.} At a fixed cap, width--depth
trades are noise. If the cap is the constraint, the honest experiment is to raise it.
\item \textbf{Verify function-preserving transforms with logits, automatically, inside the tool
that performs them.} Our graft tool eventually refused to write when
$|\Delta\text{loss}|$ exceeded a threshold. It should have compared logits from the start, and
it should have existed before the first graft rather than after the third.
\end{enumerate}

%======================================================================
\section{Limitations}

Single seed throughout: no result here carries an error bar from repetition, only from
evaluation-set size. MLM evaluation uses 15 batches ($\approx$30k tokens, $\pm$\nats{0.05}).

The \S\ref{sec:graft} graft results are exact at the logit level and do not depend on seed. The
\S\ref{sec:depth} layer-4 result replicates across three initialisations and 40+ ablation
readings and is the strongest evidence in the report. The \S\ref{sec:anneal} $11\times$ ratio
is a single pair of runs and should be read as an order of magnitude, not a coefficient. The
replay-retention control in \S\ref{sec:worked} is one transition with four corpora at four
weights, and the 0\% condition was not deliberately designed as a control. All downstream
numbers are one fine-tune run each, and \S\ref{sec:scoreboard} argues they should not be
over-read.

The comparison to BERT-base is from published figures and is not a controlled comparison:
different data, tokenizer, training budget and evaluation harness.

%======================================================================
\appendix
\section{Reproducing the headline measurements}

{\small\begin{verbatim}
# Section 4 -- the graft ladder, end to end: builds every variant, scores it,
# and runs the logit-level function-preservation check
python tools/graft_ladder.py --device cuda

# fixed cross-model MLM evaluation (the only comparable pretraining numbers)
python tools/span_eval.py \
    --ckpt checkpoints/ckpt_v2_anneal_cosmo/last.pt \
    --baseline checkpoints/ckpt3/last.pt --device cuda

# causal layer value: single + pairwise ablation cost
python tools/layer_contrib.py --full \
    --ckpt checkpoints/ckpt_v2_anneal_cosmo/last.pt

# graft variants: --new-layer-init {stack,random,zero}, --n-layers, --d-ff
python tools/build_v2_from_ckpt3.py --n-layers 5 --d-ff 3072 \
    --new-layer-init zero --out /tmp/graft.pt

# is the reported val loss measuring what it claims?
# (--data-name is one unbroken line; wrapped here for width)
python tools/audit_eval.py --device cuda \
    --ckpt checkpoints/ckpt_v2_anneal_cosmo/last.pt \
    --data-name "cosmopedia:403451249:403451249,
                 ultrafineweb_en:256832885:256832885,
                 tinystories:73377468:73377468" \
    --prior-runs "ultrafineweb_en:167000000,tinystories:167000000"

# masking-scheme difficulty offsets
python tools/masking_difficulty.py --ckpt checkpoints/ckpt3/last.pt

# attention health -- head_sim is CENTRED; the raw null is ~0.97
python tools/attn_health.py \
    --ckpt checkpoints/ckpt_v2_anneal_cosmo/last.pt
\end{verbatim}}

Numeric outputs backing the tables are committed in
\texttt{results/paper\_measurements.json} and \texttt{results/graft\_ladder.json}.

\begin{thebibliography}{99}\small
\bibitem{bert} J.\ Devlin, M.-W.\ Chang, K.\ Lee, K.\ Toutanova.
\emph{BERT: Pre-training of Deep Bidirectional Transformers for Language Understanding.}
NAACL, 2019.
\bibitem{albert} Z.\ Lan, M.\ Chen, S.\ Goodman, K.\ Gimpel, P.\ Sharma, R.\ Soricut.
\emph{ALBERT: A Lite BERT for Self-supervised Learning of Language Representations.}
ICLR, 2020.
\bibitem{spanbert} M.\ Joshi, D.\ Chen, Y.\ Liu, D.\ S.\ Weld, L.\ Zettlemoyer, O.\ Levy.
\emph{SpanBERT: Improving Pre-training by Representing and Predicting Spans.} TACL, 2020.
\bibitem{net2net} T.\ Chen, I.\ Goodfellow, J.\ Shlens.
\emph{Net2Net: Accelerating Learning via Knowledge Transfer.} ICLR, 2016.
\bibitem{stacking} L.\ Gong, D.\ He, Z.\ Li, T.\ Qin, L.\ Wang, T.-Y.\ Liu.
\emph{Efficient Training of BERT by Progressively Stacking.} ICML, 2019.
\bibitem{bert2bert} C.\ Chen et al.
\emph{bert2BERT: Towards Reusable Pretrained Language Models.} ACL, 2022.
\bibitem{data2vec} A.\ Baevski, W.-N.\ Hsu, Q.\ Xu, A.\ Babu, J.\ Gu, M.\ Auli.
\emph{data2vec: A General Framework for Self-supervised Learning in Speech, Vision and
Language.} ICML, 2022.
\bibitem{kaplan} J.\ Kaplan et al. \emph{Scaling Laws for Neural Language Models.}
arXiv:2001.08361, 2020.
\bibitem{chinchilla} J.\ Hoffmann et al.
\emph{Training Compute-Optimal Large Language Models.} arXiv:2203.15556, 2022.
\bibitem{induction} C.\ Olsson et al.
\emph{In-context Learning and Induction Heads.} Transformer Circuits Thread, 2022.
\bibitem{sbert} N.\ Reimers, I.\ Gurevych.
\emph{Sentence-BERT: Sentence Embeddings using Siamese BERT-Networks.} EMNLP, 2019.
\bibitem{tinystories} R.\ Eldan, Y.\ Li.
\emph{TinyStories: How Small Can Language Models Be and Still Speak Coherent English?}
arXiv:2305.07759, 2023.
\bibitem{cosmopedia} L.\ Ben Allal, A.\ Lozhkov, G.\ Penedo, T.\ Wolf, L.\ von Werra.
\emph{Cosmopedia.} Hugging Face, 2024.
\bibitem{minicpm} S.\ Hu et al.
\emph{MiniCPM: Unveiling the Potential of Small Language Models with Scalable Training
Strategies.} arXiv:2404.06395, 2024. (Warmup--Stable--Decay schedule.)
\end{thebibliography}

\end{document}
